\section{Conclusion}
In this study, we proposed a deep learning-based self-supervised learning restoration method that leverages SAHS images acquired via a THz-TDS system to improve image quality. This innovative learning approach does not require additional training datasets or artificial image degradation; instead, it directly utilizes the frequency characteristics of THz SAHS images. Low-frequency images, which are constrained by the diffraction limits of their frequencies, typically cannot achieve the sharpness of high-frequency images. However, our method successfully emulates high-frequency images, thereby overcoming the diffraction limits. Additionally, in typical THz systems, the SNR decreases with increasing frequency, thus resulting in noisy images at higher frequencies despite their enhanced sharpness. To address this issue, we used THz SAHS images to train models that learned low-noise images at lower frequencies, thus effectively reducing noise in high-frequency images.

Additionally, we introduce three frameworks based on self-supervised learning using THz SAHS images. First, we propose an O2O-mapping learning method to learn the mapping between two frequencies. This method enhances the sharpness of low-frequency images beyond their diffraction limits and reduces noise in high-frequency areas with a low SNR. The frequency relationship between the source, target, and inference images critically affects the inference performance. Thus, the source and target frequencies as well as the frequency of the inference image frequency must be selected appropriately to optimize the image quality. Furthermore, we demonstrated quality improvements in new THz images not used during training, which signified the broad applicability and potential of the proposed method in various THz imaging systems. Second, we propose an RE learning method that recurrently applies O2O learning to further improve the image quality. As the number of RE iterations increased, the distinctness between the index bars of the target objects in the test images became more pronounced. However, excessive iterations can lead to image distortion, thus highlighting the importance of judiciously determining the number of RE learning applications required for optimal image quality. Finally, we integrated ZSSR into our learning process to enhance image resolution. This method utilizes purely natural frequency relationships in THz SAHS images to achieve SR while overcoming diffraction limits. Integrating ZSSR into the proposed self-supervised learning process significantly enhanced image quality compared with other state-of-the-art image restoration methods.

The main contribution of this study is the proposal of a novel framework for improving image quality based on self-supervised learning using THz SAHS images. This methodology has proven to be effective in significantly enhancing the sharpness of THz images in the low-frequency domain and reducing noise in the high-frequency domain. Our method is applicable to both trained and new images, thus indicating its potential for application in various THz imaging systems. Additionally, by leveraging pure frequency relationships, our proposed learning approach is expected to improve image quality when applied across a wide range of imaging systems.